\section{Concluding remarks}
\label{sec:conclude}

In this paper, we considered asymptotic properties of the SVM in HDLSS settings
under a spiked model. We showed that the geometric representation of HDLSS data
does not hold under the spiked model, and that the HDLSS data converge to a random
configuration in a finite-dimensional space. We showed that the SVM does not hold
the consistency property, and that the bias term of the BC-SVM should be modified.
In order to overcome such difficulties, we proposed the SC-SVM and showed that it
holds the consistency property when the sample size goes to infinity. We also
showed that the growth of the sample size is essential for the projection-based
procedures. Finally, we verified Theorems~\ref{thm:1} and~\ref{thm:4},
Corollary~\ref{cor:1} and Proposition~\ref{prop:1} numerically, and we compared the
classifiers.

We give the following comparison between the existing theory and the results in
this paper.

\begin{table}[ht]
\centering
\caption{Comparison between the non-spiked model and the spiked model.}
\label{tab:model-compare}
\begin{tabular}{@{}p{3.2cm}p{5.2cm}p{5.2cm}@{}}
\toprule
Feature & non-spiked model & spiked model \\
\midrule
eigenvalue condition
& $\tr(\Sigma_i^2)/\theta_i^2\to 0$
& $\lambda_{is}/\theta\to c_{is}\in(0,1)$ \\
data configuration
& deterministic regular simplex
& random configuration in $\R^{1+m_1+m_2}$ \\
Gram matrix
& deterministic limit
& random $G^*$ \\
positive definiteness
& from the geometric representation
& only from the diagonal $D$ \\
SVM solution
& deterministic, closed form
& random limit \\
noise of $x_0$
& contributes as a diagonal term
& does not contribute \\
misclassification rate
& tends to zero by the BC-SVM
& does not tend to zero \\
bias term
& $\kappa$ in~\eqref{eq:1.3}
& $\kappa_*$ in Proposition~\ref{prop:1} \\
remedy
& bias correction
& spike removal and bias correction \\
required sample size
& $N$ may be fixed
& $n_i\to\infty$ is essential \\
\bottomrule
\end{tabular}
\end{table}

We close the paper with some remaining problems.

\begin{enumerate}[label=(\roman*)]
\item The theory of the SVM with the Gaussian kernel given by \citet{nakayama2020}
relies on the fact that the argument of the kernel function becomes deterministic
in the limit. From Theorem~\ref{thm:1}, the argument is random under the spiked
model, so that the kernel matrix itself has a random limit. It is a natural
problem to investigate how the choice of the scale parameter should be modified.

\item As we remarked in Remark~\ref{rem:7}, a minimax type lower bound for fixed
$N$ remains open.

\item The SC-SVM uses the NR methodology as a preprocessing. It is an interesting
problem to construct a framework in which the projection and the separating
hyperplane are estimated simultaneously. We also note from
Section~\ref{sec:sim-classifiers} that the sample splitting assumed in
Theorem~\ref{thm:5} is expensive in practice, so that a theory without the
splitting is desirable.

\item We assumed~(S), in which the spiked eigenvalues are of the same order as
$\theta$. On the other hand, there is an intermediate region in which
$\lambda_{i1}/\theta\to 0$ while $\lambda_{i1}^2/\tr(\Sigma_i^2)$ does not tend to
zero. In this region, the geometric representation is preserved but the second
order asymptotics is different, so that we expect that the consistency property is
preserved while the convergence rate is changed.

\item We considered the hard-margin SVM. When a regularization parameter $C$ is
introduced, the feasible set of the dual problem becomes
$\{\alpha:0\le\alpha_p\le C\}$ and the a priori bound in Lemma~\ref{lem:5} becomes
easier, so that the arguments in Section~\ref{sec:svm} carry over. On the other
hand, we expect that the limit depends on the relative scale of $C$ and $\theta$,
and the classification of the cases is not yet organized.

\item In the simulations we took the spiked direction orthogonal to $\mu$. As we
noted in Section~\ref{sec:sim-classifiers}, the advantage of the SC-SVM over the
bias correction alone is expected to be larger when the spike is not orthogonal to
$\mu$. A systematic study of this case, together with an application to actual
data, is left for future work.
\end{enumerate}
